\documentclass[11pt]{article}
\usepackage[margin=1in]{geometry}
\usepackage{amsmath,amssymb,amsthm}
\usepackage{booktabs}
\usepackage{hyperref}
\usepackage{tikz}
\usetikzlibrary{arrows.meta}

\newtheorem{theorem}{Theorem}
\newtheorem{lemma}[theorem]{Lemma}
\newtheorem{definition}[theorem]{Definition}
\newtheorem{remark}[theorem]{Remark}

\title{Bounded Perturbation Invariance for Deterministic Argmax}
\author{Research Architecture Runtime \\ \texttt{feature/lean-gated-operator-papers}}
\date{2026-07-25}

\begin{document}
\maketitle

\begin{abstract}
We study the finite deterministic argmax operator
\(\operatorname*{arg\,max}_{i\in[m]} s_i\) under additive score perturbations
bounded in the \(\ell_\infty\) norm.
When a unique maximizer \(i^\star\) exists with margin
\(\gamma(s)=s_{i^\star}-\max_{j\neq i^\star}s_j\), we prove that
\(\gamma(s)>2\varepsilon\) is necessary and sufficient for \(i^\star\) to remain
the unique maximizer under every perturbation \(\delta\) with
\(\|\delta\|_\infty\le\varepsilon\).
Repository publication is gated on Lean kernel check
(\texttt{LEAN\_FULL}) of the Mathlib $\mathbb{R}$ core (\texttt{REAL\_MATHLIB}) in
\texttt{Research.Operators.Argmax.Margin}
(\texttt{margin\_invariance}, \texttt{margin\_sharpness}).
\end{abstract}

\section{Introduction}
Selection operators such as argmax are unstable near ties:
infinitesimal score noise can change the reported index.
A quantitative margin condition restores invariance under bounded noise.
This note records the classical \(\ell_\infty\) margin theorem as the first
operator-specific result exercised end-to-end through the Discovery and
Verification architectures in this repository.

\section{Problem formulation}
Fix \(m\ge 2\) and scores \(s\in\mathbb{R}^m\).
Write \([m]=\{1,\dots,m\}\).
The (set-valued) argmax is
\[
\operatorname*{Argmax}(s)=\bigl\{i\in[m]: s_i=\max_{j\in[m]} s_j\bigr\}.
\]
We write \(i^\star=\operatorname*{arg\,max}_i s_i\) when
\(\lvert\operatorname*{Argmax}(s)\rvert=1\).

\paragraph{Perturbation model.}
For \(\varepsilon\ge 0\), admissible perturbations satisfy
\(\|\delta\|_\infty=\max_i\lvert\delta_i\rvert\le\varepsilon\).
The perturbed scores are \(s+\delta\).

\section{Notation and definitions}
\begin{definition}[Margin]
If \(i^\star\) is the unique maximizer of \(s\), the margin is
\[
\gamma(s)=s_{i^\star}-\max_{j\neq i^\star}s_j.
\]
Otherwise \(\gamma(s)\) is undefined.
\end{definition}

\section{Mathematical assumptions}
\begin{enumerate}
\item \(m\ge 2\) and \(s\in\mathbb{R}^m\).
\item A unique maximizer \(i^\star\) exists, so \(\gamma(s)>0\).
\item Perturbations obey \(\|\delta\|_\infty\le\varepsilon\) with \(\varepsilon\ge 0\).
\end{enumerate}
We do not treat data-dependent scores \(s_i(D)\), other \(\ell_p\) balls, or
set-valued selection rules beyond recording that ties exclude the hypothesis.

\section{Main theorem}
\begin{theorem}[Bounded perturbation invariance]\label{thm:main}
Under the assumptions above, if \(\gamma(s)>2\varepsilon\), then for every
admissible \(\delta\), \(i^\star\) is the unique maximizer of \(s+\delta\).
\end{theorem}

\begin{theorem}[Sharpness]\label{thm:sharp}
Under the same assumptions, if \(\gamma(s)\le 2\varepsilon\), then there exists
an admissible \(\delta\) such that \(i^\star\) is not the unique maximizer of
\(s+\delta\).
\end{theorem}

\section{Supporting lemmas}
\begin{lemma}[Worst-case gap]\label{lem:gap}
For any admissible \(\delta\) and any \(j\neq i^\star\),
\[
(s+\delta)_{i^\star}-(s+\delta)_j
\ge s_{i^\star}-s_j-2\varepsilon
\ge \gamma(s)-2\varepsilon.
\]
\end{lemma}

\begin{proof}
\((s+\delta)_{i^\star}\ge s_{i^\star}-\varepsilon\) and
\((s+\delta)_j\le s_j+\varepsilon\).
Subtracting yields the first inequality; the second uses the definition of
\(\gamma\).
\end{proof}

\section{Proofs}
\begin{proof}[Proof of Theorem~\ref{thm:main}]
Lemma~\ref{lem:gap} and \(\gamma(s)>2\varepsilon\) imply
\((s+\delta)_{i^\star}>(s+\delta)_j\) for all \(j\neq i^\star\).
\end{proof}

\begin{proof}[Proof of Theorem~\ref{thm:sharp}]
Let \(j^\star\in\operatorname*{Argmax}_{j\neq i^\star}s_j\).
Set \(\delta_{i^\star}=-\varepsilon\), \(\delta_{j^\star}=+\varepsilon\), and
\(\delta_k=0\) otherwise.
Then \(\|\delta\|_\infty\le\varepsilon\) and
\[
(s+\delta)_{i^\star}-(s+\delta)_{j^\star}=\gamma(s)-2\varepsilon\le 0,
\]
so \(i^\star\) is not a unique maximizer.
\end{proof}

\section{Discussion}
\subsection{Intuition}
The margin is a buffer against adversarial noise that simultaneously depresses
the winner and elevates a runner-up. Each direction can consume at most
\(\varepsilon\), hence the factor two.

\subsection{Geometric explanation}
\begin{figure}[h]
\centering
\begin{tikzpicture}[scale=1.1]
\draw[->] (0,0) -- (5.2,0) node[right] {\(s\)};
\draw[thick] (0,1.2) -- (1.6,1.2) node[above left] {\(s_{j^\star}\)};
\draw[thick] (0,3.0) -- (3.2,3.0) node[above left] {\(s_{i^\star}\)};
\draw[<->] (3.4,1.2) -- (3.4,3.0) node[midway,right] {\(\gamma\)};
\draw[dashed] (0,3.0-0.7) -- (4.5,3.0-0.7);
\draw[dashed] (0,1.2+0.7) -- (4.5,1.2+0.7);
\node at (4.7,2.65) {\small \(-\varepsilon\)};
\node at (4.7,1.75) {\small \(+\varepsilon\)};
\node[align=center,font=\small] at (2.2,-0.7)
{Invariance requires \(\gamma>2\varepsilon\)};
\end{tikzpicture}
\caption{Margin buffer versus \(\ell_\infty\) adversarial moves.}
\label{fig:margin}
\end{figure}

\subsection{Why the theorem matters}
It gives an exact, checkable certificate that a reported argmax is stable to
bounded score noise---useful for post-hoc selection and audit of ranking
pipelines.

\subsection{Where assumptions are used}
Uniqueness supplies \(\gamma\); \(\ell_\infty\) yields the additive \(2\varepsilon\)
shrinkage; finiteness of \(m\) is used only to ensure maxima exist.
Indices in the paper are \([m]=\{1,\dots,m\}\); the reference implementation
uses \(0..m-1\) with the obvious shift.

\subsection{Examples}
Take \(s=(3,1,0.5)\) so \(\gamma=2\). For \(\varepsilon=0.9\), invariance holds.
For \(\varepsilon=1\), \(\delta=(-1,+1,0)\) yields \(s+\delta=(2,2,0.5)\) (tie).

\subsection{Counterexamples and boundary cases}
At \(\gamma=2\varepsilon\) the adversary forces a tie: uniqueness fails though
\(i^\star\) may remain \emph{a} maximizer.
Ties in \(s\) exclude the theorem hypothesis entirely.

\section{Interpretation}
The condition \(\gamma>2\varepsilon\) is both a proof target and an operational
test: compute the empirical margin and compare to the noise budget.

\section{Utility implications}
Stabilization mechanisms (score smoothing, hysteresis, top-\(k\) reporting) can
be justified when the measured margin is too small relative to \(\varepsilon\).

\section{Limitations}
\begin{itemize}
\item Lean formalization is a Mathlib \(\mathbb{R}\) core
  (\texttt{REAL\_MATHLIB}) scoped to the margin theorems.
\item Finite score vectors; abstract \(s\), not \(s_i(D)\).
\item \(\ell_\infty\) only; other norms change constants.
\item Engineering float checks can blur exact equality \(\gamma=2\varepsilon\); Lean uses exact \(\mathbb{R}\).
\end{itemize}

\section{Open questions}
Extensions to ties (set-valued argmax), \(\ell_p\) balls, and data-dependent
scores \(s_i(D)\); composition with top-\(k\) and thresholding.

\section{Connections to existing framework}
Artifacts are minted as typed Discovery IR, packaged as a sealed CRP, System~B
discharges obligations, then System~3 runs \texttt{lake build} + axiom capture.
Lean modules: \texttt{Research.Operators.Argmax.Margin}.
Certificate: \texttt{lean/certificates/argmax/bounded-perturbation-margin/}.

\section{Verifier summary}
Publication requires \texttt{derived\_lean\_status}=\texttt{LEAN\_FULL}
(sorry/admit \(=0\), \texttt{build\_ok}, axiom closure captured).
Kernel axioms observed: \texttt{propext}, \texttt{Quot.sound}.
Prior computational-only paper archived under
\texttt{\_archive/v0-computational/}.

\section{Appendix}
Canonical statement (repository string):
\begin{quote}\small
Let \(m\ge 2\), \(s\in\mathbb{R}^m\) with unique maximizer
\(i^\star=\operatorname*{arg\,max}_i s_i\) and margin
\(\gamma(s)=s_{i^\star}-\max_{j\neq i^\star}s_j>0\).
Let \(\varepsilon\ge 0\) and \(\delta\in\mathbb{R}^m\) with \(\|\delta\|_\infty\le\varepsilon\).
If \(\gamma(s)>2\varepsilon\), then \(i^\star\) is the unique maximizer of \(s+\delta\).
\end{quote}

\section*{References}
\begin{thebibliography}{9}
\bibitem{repo}
Repository implementation:
\texttt{implementation/src/operators/argmax/}
and \texttt{implementation/docs/operators/argmax.md}.
\end{thebibliography}

\end{document}
